\documentclass[10pt]{article}
\usepackage[margin=1in]{geometry}
\usepackage[T1]{fontenc}
\usepackage{lmodern,microtype,amsmath,amssymb,amsthm}
\theoremstyle{definition}
\newtheorem{assumption}{Assumption}
\begin{document}
\section*{Gradient Descent for Gaussian Recovery}
We study recovery of two unknown means $\mu_0,\mu_1\in\mathbb{R}^d$, where $d\geq1$, from exact gradients of a reverse Fisher loss for equal-weight Gaussian mixtures with identity covariance. Their separation lies in $[a,1]$, where $a\in[0,1]$, and their midpoint has norm at most $C\geq0$. The loss $F_\mu$ averages the squared difference between candidate and target scores under the candidate mixture. The oracle returns $\nabla F_\mu(\theta)$ at $\theta\in(\mathbb{R}^d)^2$. Constant-step gradient descent must return two means whose summed squared distances to the target means, minimized over the two matchings, are at most $\varepsilon\in(0,1]$ with probability at least $1-\delta$, where $\delta\in(0,1]$.

Write $E=\mathbb{R}^d$ with its Euclidean norm and equip $E^2$ with the product norm
\[
\|\theta\|^2=\|\theta_0\|^2+\|\theta_1\|^2.
\]
For a candidate pair $\theta=(\theta_0,\theta_1)\in E^2$, define its mixture density and score by
\[
\begin{aligned}
q_\theta(x)
&=\frac{1}{2(2\pi)^{d/2}}\sum_{i=0}^1
\exp\!\left(-\frac{\|x-\theta_i\|^2}{2}\right),\\
w_i(\theta,x)
&=\frac{\exp(-\|x-\theta_i\|^2/2)}
{\sum_{j=0}^1\exp(-\|x-\theta_j\|^2/2)},\\
s_\theta(x)
&=\nabla_x\log q_\theta(x)
=\sum_{i=0}^1 w_i(\theta,x)\theta_i-x.
\end{aligned}
\]
The reverse Fisher loss and the associated optimization problem are
\[
F_\mu(\theta)
=\frac12\sum_{i=0}^1
\mathbb{E}_{Z\sim\mathcal{N}(0,I_d)}
\left[\|s_\theta(\theta_i+Z)-s_\mu(\theta_i+Z)\|^2\right],
\qquad
\min_{\theta\in E^2}F_\mu(\theta).
\]
Here $I_d$ is the identity matrix. Both candidate means are free optimization variables. The algorithm receives $d,C,a,\varepsilon,\delta$; the target pair $\mu$ is unknown.

Our assumptions:
\begin{assumption}[Target geometry]
The target means satisfy
\[
\begin{aligned}
a&\leq\|\mu_0-\mu_1\|\leq1,\\
\left\|\frac{\mu_0+\mu_1}{2}\right\|&\leq C.
\end{aligned}
\]
\end{assumption}

Queries may be made at any $\theta\in E^2$. The oracle returns $\nabla F_\mu(\theta)$ at $\theta$, with the gradient taken in the product Euclidean space.

Our result is algorithm-specific: we analyze constant-step gradient descent. The algorithm chooses an initialization $\theta^0\in E^2$, which may be random, and a fixed step size $\gamma>0$. Each query is at the current iterate, and the update is
\[
\theta^{k+1}=\theta^k-\gamma\nabla F_\mu(\theta^k).
\]
After $K\in\mathbb{N}_0$ oracle calls, the output is the last iterate $\bar{\theta}=\theta^K$. The initialization and step size are chosen using the available inputs.

Define the squared matching error by
\[
\begin{aligned}
M_\mu(\theta)=\min\bigl\{
&\|\theta_0-\mu_0\|^2+\|\theta_1-\mu_1\|^2,\\
&\|\theta_0-\mu_1\|^2+\|\theta_1-\mu_0\|^2
\bigr\}.
\end{aligned}
\]
The output must be measurable and satisfy
\[
P\bigl(M_\mu(\bar{\theta})\leq\varepsilon\bigr)\geq1-\delta,
\]
where $P$ is the law of the algorithm's randomness.
\end{document}
